\documentclass[12pt,a4paper]{article}
% Globalni nastaveni dokumentu.
\usepackage[a4paper,left=3.2cm,right=3.2cm,top=2.7cm,bottom=2.9cm,headheight=15pt]{geometry}
\usepackage{fontspec}
\usepackage{polyglossia}
\setdefaultlanguage{czech}
\setmainfont{TeX Gyre Pagella}
\setsansfont{TeX Gyre Heros}
\usepackage{microtype}
\usepackage{enumitem}
\usepackage{hyperref}
\usepackage{xcolor}
\usepackage{amsmath,amssymb}
\DeclareMathOperator*{\plim}{plim}
\usepackage{titlesec}
\usepackage{tocloft}
\setlength{\cftsecnumwidth}{2.6em}
\usepackage{fancyhdr}
\hypersetup{colorlinks=true,linkcolor=black,urlcolor=blue}
\linespread{1.08}
\setlength{\parindent}{0pt}
\setlength{\parskip}{0.45\baselineskip plus 1pt}
\setlength{\emergencystretch}{2em}

\makeatletter
\renewcommand*\l@section{\@dottedtocline{1}{0em}{2.4em}}
\makeatother

\setlist[itemize]{topsep=4pt,itemsep=2pt,parsep=0pt,leftmargin=1.4em}
\setlist[enumerate]{topsep=4pt,itemsep=2pt,parsep=0pt,leftmargin=1.7em}
\titleformat{\section}{\Large\bfseries\sffamily}{\thesection}{0.6em}{}
\titleformat{\subsection}{\large\bfseries\sffamily}{\thesubsection}{0.6em}{}
\titleformat{\subsubsection}{\normalsize\bfseries\sffamily}{\thesubsubsection}{0.6em}{}
\titlespacing*{\section}{0pt}{1.2\baselineskip}{0.55\baselineskip}
\titlespacing*{\subsection}{0pt}{0.9\baselineskip}{0.35\baselineskip}
\titlespacing*{\subsubsection}{0pt}{0.7\baselineskip}{0.25\baselineskip}
\pagestyle{fancy}
\fancyhf{}
\lhead{\small Příprava ke státnicím - Data Analytics}
\rhead{\small\thepage}
\cfoot{}
\newcommand{\term}[1]{\textbf{#1}}
\newcommand{\uv}[1]{„#1“}
\newcommand{\todohead}{\subsection*{Co doplnit mimo dodané materiály}}
\newcommand{\pozor}{\subsection*{Pozor u zkoušky}}


\lhead{\small Posudkové otázky - Data Analytics}

\begin{document}

\begin{titlepage}
  \centering
  \vspace*{2.4cm}
  {\Large\bfseries Bakalářské studium\par}
  \vspace{1.2cm}
  {\Huge\bfseries Posudkové otázky ke státnicím\par}
  \vspace{0.6cm}
  {\LARGE\bfseries Data Analytics\par}
  \vspace{0.8cm}
  \rule{0.72\textwidth}{0.4pt}\par
  \vspace{0.8cm}
  {\large Vybrané otázky z posudků a odborných rozprav, seskupené podle uvedené komise\par}
  \vfill
\end{titlepage}

\tableofcontents
\newpage

\section{Holý, Černý, Chudán}
Uvedená komise: doc. Mgr. Vladimír Holý, Ph.D.; PhDr. Jan Černý, Ph.D.; Ing. David Chudán, Ph.D.

\begin{enumerate}
  \item \term{Jaké problémy vznikají, když regresní model pracuje s časovou řadou?}

  Časové řady často porušují předpoklad nezávislých pozorování. Typické problémy jsou trend,
  sezónnost, autokorelace chyb, nestacionarita, strukturální zlomy a proměnlivý rozptyl v čase.
  Při modelování je potřeba testovat stacionaritu, pracovat se zpožděními, sezónními dummy proměnnými
  nebo ARIMA/SARIMA strukturou a používat vhodné standardní chyby.

  \item \term{Co je zdánlivá regrese?}

  Zdánlivá regrese vzniká typicky u nestacionárních časových řad, které mají společný trend, ale
  věcně spolu nemusí souviset. Model může mít vysoké $R^2$ a významné koeficienty, přestože vztah je
  falešný. Řešením je zkoumat stacionaritu, pracovat s diferencemi, ověřit kointegraci nebo modelovat
  vztah v podobě vhodné pro časové řady.

  \item \term{Jak řešit autokorelovanou náhodnou složku?}

  Nejprve je potřeba autokorelaci diagnostikovat, například reziduálním grafem, ACF/PACF nebo
  Durbinovým-Watsonovým testem. Možné reakce jsou doplnit zpožděné proměnné, modelovat chyby jako
  AR proces, použít GLS nebo použít robustní standardní chyby pro časové řady, například Newey-West.
  V predikčních úlohách může dávat smysl přejít na modely ARIMA/SARIMA nebo dynamickou regresi.

  \item \term{Co je etika a informační etika?}

  Etika je systematické uvažování o tom, co je správné jednání. V datové analytice nejde jen o právní
  soulad, ale také o přiměřenost, férovost, transparentnost a odpovědnost. Informační etika se týká
  sběru, uchovávání, sdílení, zpracování a používání informací. Řeší soukromí, přesnost, vlastnictví
  informací, dostupnost, odpovědnost za rozhodnutí a riziko poškození lidí chybným nebo neférovým
  výstupem.

  \item \term{Co jsou otevřená data a jak se liší od veřejně dostupných dat?}

  Otevřená data jsou zveřejněná tak, aby byla strojově čitelná, dobře dokumentovaná, licenčně
  použitelná a ideálně dostupná přes hromadné stažení nebo API. Veřejně dostupná data mohou být
  pouze zobrazena na webu, ale nemusí mít licenci pro další použití, vhodný formát ani metadata.
  PDF na webu tedy může být veřejné, ale nemusí být skutečně otevřeným datovým zdrojem.

  \item \term{Jaké je základní členění strojového učení a jaké znáte unsupervised metody?}

  Základní členění je učení s učitelem, učení bez učitele, učení s posilováním a někdy také
  semi-supervised/self-supervised učení. Učení bez učitele hledá strukturu bez cílové proměnné.
  Typické metody jsou shlukování, PCA a jiné redukce dimenze, asociační pravidla, detekce anomálií
  a topic modeling.

  \item \term{Jak fungují rozhodovací stromy a jaké existují další klasifikační metody?}

  Rozhodovací strom rekurzivně dělí data podle pravidel typu \(x_j \le c\), aby v listech vznikly
  co nejčistší skupiny. Pro klasifikaci se používá například Gini impurity nebo informační zisk.
  Další klasifikační metody jsou logistická regrese, k-NN, naive Bayes, SVM, rozhodovací lesy,
  gradient boosting a neuronové sítě.

  \item \term{Co znamená \uv{random} v metodě random forest?}

  Random forest kombinuje mnoho rozhodovacích stromů. Náhodnost vzniká jednak bootstrapováním
  pozorování pro každý strom, jednak náhodným výběrem podmnožiny prediktorů při hledání splitu.
  Cílem je snížit korelaci mezi stromy a tím snížit rozptyl výsledného modelu.

  \item \term{Jaký je rozdíl mezi Business Intelligence a Competitive Intelligence?}

  Business Intelligence se zaměřuje hlavně na interní data organizace a jejich přeměnu na reporty,
  dashboardy, KPI a rozhodovací podporu. Competitive Intelligence sleduje vnější prostředí:
  konkurenci, trh, zákazníky, technologie, regulaci a veřejné zdroje. BI typicky integruje provozní
  data do datového skladu, CI častěji pracuje s externími zdroji, OSINT, webem, registry a tržními
  informacemi.

  \item \term{Jak se předzpracovávají nestrukturovaná data a jaké jsou typické úlohy NLP?}

  Text se obvykle čistí, normalizuje, tokenizuje, převádí na lemmata nebo stemmy a reprezentuje pomocí
  bag-of-words, TF-IDF, embeddingů nebo jazykových modelů. U dokumentů se řeší také indexace a
  vyhledávání. Typické úlohy NLP jsou klasifikace textů, sentiment analysis, extrakce entit,
  topic modeling, sumarizace, překlad, informační vyhledávání a question answering.

  \item \term{Jak modelovat sezónnost v regresních modelech a časových řadách?}

  V regresi lze použít sezónní dummy proměnné, Fourierovy členy, trend, zpožděné hodnoty a interakce.
  V časových řadách se používají sezónní diference a sezónní AR/MA složky. SARIMA model se značí
  \((p,d,q)(P,D,Q)_s\), kde \(s\) je délka sezóny, například 4 pro čtvrtletní data nebo 12 pro
  měsíční data.

  \item \term{Co znamená exploratorní analýza dat?}

  Exploratorní analýza dat je fáze poznávání dat před modelováním. Hledají se chyby, chybějící
  hodnoty, odlehlá pozorování, rozdělení proměnných, vztahy mezi proměnnými a možné hypotézy.
  Používají se deskriptivní statistiky, histogramy, boxploty, scatterploty, korelace a datové
  profilování. EDA nemá dokazovat definitivní závěry, ale pomáhá dobře formulovat další analýzu.

  \item \term{Co je informační chování a Wilsonův model?}

  Informační chování popisuje, jak lidé rozpoznají informační potřebu, hledají informace, vybírají
  zdroje, vyhodnocují relevanci a informace používají. Wilsonův model zdůrazňuje kontext uživatele,
  bariéry hledání, aktivaci informační potřeby a způsoby vyhledávání. Pro státnice stačí říct, že
  vyhledávání není jen technický dotaz do systému, ale lidský proces ovlivněný potřebou, situací,
  motivací a dostupností zdrojů.

  \item \term{Co je booleovský model vyhledávání?}

  Booleovský model reprezentuje dokumenty a dotaz pomocí termů a logických operátorů AND, OR a NOT.
  Dokument buď podmínku splní, nebo nesplní; klasický model tedy nepracuje s plynulým skóre relevance.
  Výhodou je přesná kontrola dotazu, nevýhodou křehkost: příliš úzký dotaz může vrátit málo výsledků,
  příliš široký dotaz mnoho nerelevantních dokumentů.

  \item \term{Jaké znáte heuristické a metaheuristické optimalizační algoritmy?}

  Heuristiky hledají dobré řešení bez záruky globální optimality. Patří sem greedy algoritmy,
  lokální hledání, hill climbing, tabu search, simulated annealing, genetické algoritmy, ant colony
  optimization a particle swarm optimization. Používají se u rozsáhlých kombinatorických úloh, kde
  přesné metody jako úplné prohledávání nebo MILP mohou být příliš pomalé.

  \item \term{Co je OSINT a jak souvisí s externími zdroji?}

  OSINT je získávání a vyhodnocování informací z otevřených zdrojů: webů, registrů, médií, sociálních
  sítí, akademických databází, veřejných dokumentů a otevřených dat. V analytice slouží k doplnění
  interních dat o kontext trhu, konkurence, událostí nebo rizik. I u OSINT je nutné řešit kvalitu,
  licenci, metodiku, etiku a ochranu osobních údajů.

  \item \term{Jak se vyhodnocují metody strojového učení?}

  U klasifikace se používá confusion matrix, accuracy, precision, recall, F1, ROC křivka a AUC.
  U regrese například MAE, MSE, RMSE a $R^2$. Důležité je oddělit trénovací, validační a testovací
  data, případně použít k-fold cross-validation. Metrika se volí podle cíle: jiná je vhodná pro
  nevyváženou klasifikaci a jiná pro symetrickou regresní chybu.

  \item \term{Jaké jsou předpoklady lineární regrese a co se stane při jejich porušení?}

  Základní předpoklady jsou linearita v parametrech, náhodný výběr nebo vhodný výběrový mechanismus,
  žádná perfektní multikolinearita, exogenita, homoskedasticita a pro malé vzorky často normalita
  chyb. Porušení exogenity vede k vychýlení, heteroskedasticita a autokorelace typicky kazí standardní
  chyby, multikolinearita zvyšuje nejistotu koeficientů a špatný funkční tvar vede k chybné
  interpretaci.

  \item \term{Co znamená korelace 0,8 mezi proměnnými?}

  Korelace 0,8 znamená silnou kladnou lineární souvislost, ale ne důkaz kauzality. V regresi mezi
  vysvětlujícími proměnnými může jít o multikolinearitu, která zvětšuje standardní chyby a zhoršuje
  interpretaci jednotlivých koeficientů. Pokud jde o korelaci mezi predikcí a skutečností, ukazuje
  silnou shodu, ale je potřeba doplnit i chybu a validaci mimo trénovací data.

  \item \term{Jaký je životní cyklus datového projektu a co se děje při nasazení?}

  Typický cyklus lze popsat přes CRISP-DM: business understanding, data understanding, data preparation,
  modeling, evaluation a deployment. Při nasazení se řeší integrace do procesu, přístupová práva,
  dokumentace, monitoring kvality dat, monitoring výkonu modelu, školení uživatelů, provozní odpovědnost
  a plán údržby. Projekt je úspěšný až tehdy, když výstup skutečně podporuje rozhodování.

  \item \term{Jak zajistit, aby se výstup datového projektu používal?}

  Je nutné zapojit uživatele už při zadání, navázat výstup na konkrétní rozhodnutí, dobře definovat
  metriky úspěchu, zajistit srozumitelnou interpretaci, školení, podporu a vlastnictví procesu.
  U dashboardů pomáhá governance metrik, kvalita dat a pravidelný refresh. U modelů je potřeba
  monitoring, vysvětlitelnost a možnost zpětné vazby.

  \item \term{Jaké jsou metody a nástroje pro vizualizaci dat?}

  Základními metodami jsou tabulky, sloupcové a spojnicové grafy, histogramy, boxploty, scatterploty,
  heatmapy, mapy, dashboardy a interaktivní reporty. Nástroje zahrnují Excel, Power BI, Tableau, Qlik,
  Python knihovny, R a webové knihovny. Volba grafu se řídí otázkou: porovnání, trend, rozdělení,
  vztah, geografie nebo struktura.

  \item \term{Jaké jsou principy big data?}

  Big data se obvykle popisují pomocí charakteristik volume, velocity, variety, veracity a value.
  Nejde jen o velký objem, ale i o rychlost vzniku, různorodost formátů, nejistou kvalitu a schopnost
  vytěžit užitečnou hodnotu. Typické nástroje jsou distribuované úložiště, Spark, cloudové datové
  platformy, streamové systémy a data lake/lakehouse.

  \item \term{Co jsou random effects v panelových datech?}

  Random effects model předpokládá, že jednotkový nepozorovaný efekt je náhodný a nekoreluje s
  vysvětlujícími proměnnými. Oproti fixed effects umožňuje odhadnout i vliv časově neměnných proměnných,
  ale vyžaduje silnější předpoklad exogenity jednotkového efektu. Pokud tento předpoklad neplatí,
  bývá vhodnější fixed effects model.
\end{enumerate}

\section{Zouhar, Chudán, Zimmermann}
Uvedená komise: doc. Ing. Jan Zouhar, Ph.D.; Ing. David Chudán, Ph.D.; Ing. Pavel Zimmermann, Ph.D.

\begin{enumerate}
  \item \term{Jmenujte zástupce učení s učitelem.}

  Učení s učitelem pracuje s cílovou proměnnou. Pro regresi lze uvést lineární regresi, rozhodovací
  stromy, random forest, gradient boosting a neuronové sítě. Pro klasifikaci logistickou regresi,
  naive Bayes, k-NN, SVM, rozhodovací stromy, ensemble metody a neuronové sítě.

  \item \term{Jaké jsou typické aplikace shlukování?}

  Shlukování se používá pro segmentaci zákazníků, hledání podobných produktů, seskupování dokumentů,
  detekci anomálií, kompresi dat, předzpracování pro další modelování a exploraci dat. Výsledkem
  není \uv{správná třída}, ale navržená struktura podobnosti, kterou je nutné věcně interpretovat.

  \item \term{Jaké jsou základní výstupy asociačních pravidel?}

  Výstupem jsou pravidla typu \(X \Rightarrow Y\), například \uv{kdo koupí X, často koupí Y}.
  Hodnotí se pomocí support, confidence, lift a dalších měr. Support říká, jak často se položky
  vyskytují, confidence podmíněnou pravděpodobnost důsledku a lift porovnání vůči náhodné nezávislosti.

  \item \term{Popište SHAP values a jejich interpretaci.}

  SHAP values vysvětlují predikci modelu jako součet příspěvků jednotlivých proměnných vůči základní
  hodnotě. Kladná SHAP hodnota predikci zvyšuje, záporná ji snižuje. Výhodou je jednotný rámec pro
  lokální vysvětlení jedné predikce i globální přehled důležitosti proměnných. Nevýhodou je výpočetní
  náročnost a riziko chybné kauzální interpretace.

  \item \term{Vysvětlete LASSO, Ridge a Elastic Net. Jak škálování ovlivní výsledky?}

  Ridge přidává \(L_2\) penalizaci a zmenšuje koeficienty, ale obvykle je nenuluje. LASSO používá
  \(L_1\) penalizaci a může koeficienty nastavit přesně na nulu, takže vybírá proměnné. Elastic Net
  kombinuje obě penalizace. Škálování je důležité, protože penalizace závisí na velikosti koeficientů;
  proměnné s různými jednotkami by byly penalizovány nerovnoměrně.

  \item \term{Jak interpretovat koeficienty v logistické regresi?}

  Logistická regrese modeluje logaritmus odds:
  \[
    \log\frac{p}{1-p}=x^\top\beta.
  \]
  Koeficient \(\beta_j\) říká změnu log-odds při jednotkové změně \(x_j\). Po exponenciaci
  \(e^{\beta_j}\) dostaneme odds ratio. Pokud je \(e^{\beta_j}=1{,}5\), odds jsou při jednotkové změně
  proměnné 1,5krát vyšší, za jinak stejných podmínek.

  \item \term{Jak volit práh klasifikace a co je ROC křivka?}

  Model často vrací skóre nebo pravděpodobnost a prahová hodnota rozhoduje, kdy případ označíme jako
  pozitivní. Nižší práh zvyšuje citlivost, ale může snížit přesnost. ROC křivka zobrazuje vztah mezi
  true positive rate a false positive rate při různých prazích. AUC shrnuje schopnost modelu řadit
  pozitivní případy před negativní.

  \item \term{Co jsou algoritmy vyhledávání dokumentů a booleovský model?}

  Vyhledávání dokumentů může být založeno na booleovském modelu, vektorovém modelu, BM25 nebo
  embeddingovém vyhledávání. Booleovský model používá operátory AND, OR a NOT a dokument splní, nebo
  nesplní dotaz. Vektorové a probabilistické modely navíc řadí dokumenty podle relevance.
\end{enumerate}

\section{Kliegr, Holý, Kučera}
Uvedená komise: doc. Ing. Tomáš Kliegr, Ph.D.; doc. Mgr. Vladimír Holý, Ph.D.; Ing. Jan Kučera, Ph.D.

\begin{enumerate}
  \item \term{Jak řešit chybějící pozorování u regresních modelů?}

  Nejprve je nutné určit mechanismus chybění: MCAR, MAR nebo MNAR. Jednoduché možnosti jsou odstranění
  neúplných řádků nebo doplnění průměrem/mediánem, ale mohou být zkreslující. Lepší bývá imputace
  podle modelu, multiple imputation, indikátor chybění nebo metody, které chybějící hodnoty zvládají
  přímo. Vždy je potřeba dopad otestovat citlivostní analýzou.

  \item \term{Co se stane, když data nejsou missing completely at random?}

  Pokud chybění souvisí s pozorovanými nebo nepozorovanými hodnotami, prosté vyřazení řádků může
  způsobit výběrové zkreslení. Odhady koeficientů i průměrů pak nemusí reprezentovat populaci. U MAR
  lze pomoci modelovou imputací podle známých proměnných, u MNAR je problém těžší a vyžaduje model
  mechanismu chybění nebo citlivostní scénáře.

  \item \term{Jak charakterizovat SQL a jaké dotazy použít pro agregace?}

  SQL je deklarativní jazyk pro práci s relačními databázemi. Základní dotaz používá \texttt{SELECT},
  \texttt{FROM}, \texttt{WHERE}, \texttt{JOIN}, \texttt{GROUP BY}, agregační funkce a \texttt{HAVING}.
  Pro agregace se používá například \texttt{COUNT}, \texttt{SUM}, \texttt{AVG}, \texttt{MIN} a
  \texttt{MAX}. \texttt{WHERE} filtruje řádky před agregací, \texttt{HAVING} filtruje skupiny po
  agregaci.

  \item \term{Co je databázová transakce a proč seskupovat SQL příkazy do transakcí?}

  Transakce je logická jednotka práce v databázi. Buď proběhne celá, nebo se změny vrátí zpět. Má
  splňovat ACID: atomicitu, konzistenci, izolaci a trvanlivost. Seskupení příkazů do transakce chrání
  databázi před neúplnými změnami, například aby se při převodu peněz neodepsala částka bez jejího
  připsání.

  \item \term{Co je one-hot encoding?}

  One-hot encoding převádí kategoriální proměnnou na sadu binárních příznaků. Například barva
  \(\{\text{červená}, \text{modrá}, \text{zelená}\}\) se převede na tři sloupce s hodnotami 0/1.
  U lineární regrese s interceptem se obvykle jedna kategorie vynechá jako referenční, aby nevznikla
  perfektní multikolinearita.

  \item \term{Co je p-hodnota?}

  P-hodnota je pravděpodobnost, že za platnosti nulové hypotézy dostaneme stejně extrémní nebo
  extrémnější testovou statistiku, než pozorujeme. Malá p-hodnota znamená, že pozorovaný výsledek je
  pod nulovou hypotézou nepravděpodobný. Není to pravděpodobnost, že nulová hypotéza je pravdivá.

  \item \term{Jaké metodiky použít pro řízení datového projektu a co je DMBOK?}

  Pro analytické projekty se hodí CRISP-DM, SEMMA, TDSP nebo agilní řízení s inkrementy. Pro obecné
  projektové řízení lze zmínit PMBOK, PRINCE2 nebo Scrum. DMBOK je znalostní rámec pro data management,
  ne metodika jednoho projektu. Řeší data governance, data quality, metadata, datovou architekturu,
  bezpečnost a správu kmenových dat.

  \item \term{Jaké metody použít pro binární predikci typu obsazeno/neobsazeno?}

  Jde o binární klasifikaci. Použít lze logistickou regresi, rozhodovací strom, random forest, gradient
  boosting, SVM, naive Bayes nebo neuronovou síť. Výběr závisí na velikosti dat, interpretovatelnosti,
  nelinearitách a nákladech chyb. Vyhodnocení by mělo zahrnovat confusion matrix, precision, recall,
  F1, ROC/AUC a vhodný práh.

  \item \term{Jaké metody lze použít pro krátké panely a regionální/obecní data?}

  U panelových dat lze použít pooled OLS, fixed effects, random effects, first differences nebo
  dynamické panelové modely podle délky panelu a cíle analýzy. U krátkých panelů je potřeba hlídat
  malý počet časových období, heterogenitu jednotek, časové efekty, prostorovou závislost a robustní
  standardní chyby.

  \item \term{Jak se jmenuje první uzel rozhodovacího stromu a co znamená?}

  První uzel je kořenový uzel. Obsahuje celý trénovací soubor a první rozhodovací pravidlo, které
  model vybral jako nejpřínosnější pro rozdělení dat. Pro interpretaci je důležitý, protože ukazuje
  proměnnou a hranici, která v daném stromu nejvíce snižuje nečistotu podle zvoleného kritéria.

  \item \term{Jak postupovat při ověření hypotézy o negativních příspěvcích na sociálních sítích?}

  Nejprve definovat hypotézu, populaci příspěvků, období, zdroje a jednotku analýzy. Poté sesbírat
  data, vyčistit text, vytvořit nebo použít sentiment model a získat skóre sentimentu. Následně lze
  testovat podíl negativních příspěvků, porovnat jej s běžným obdobím nebo benchmarkem a vyhodnotit
  nejistotu. Dataset by měl obsahovat text, čas, zdroj, metadata, případně ruční label sentimentu pro
  validaci.

  \item \term{Jak se rozhodovací strom učí z dat a jaké algoritmy znáte?}

  Strom se učí rekurzivním dělením dat. V každém uzlu hledá proměnnou a hranici, která nejvíce zlepší
  čistotu skupin. Pro klasifikaci se používá Gini impurity nebo entropy, pro regresi snížení rozptylu
  či MSE. Známé algoritmy jsou ID3, C4.5, CART a jejich moderní implementace v ensemble metodách.

  \item \term{Jaké jsou právní a etické otázky web scrapingu?}

  Je potřeba respektovat autorská a databázová práva, smluvní podmínky webu, ochranu osobních údajů,
  přiměřené zatížení serveru a bezpečnostní omezení. Scraping může být problematický, pokud obchází
  technické zábrany, stahuje neveřejná data, porušuje licenci, zpracovává osobní údaje bez právního
  titulu nebo narušuje provoz služby.

  \item \term{Jak robot pozná, které stránky může stáhnout?}

  Základní vodítko je soubor \texttt{robots.txt}, který uvádí pravidla pro roboty podle user-agent.
  Může zakazovat určité cesty nebo uvádět mapu webu. Právně to samo o sobě nemusí vyřešit všechny
  otázky, ale z hlediska etikety web crawlerů je vhodné pravidla respektovat, používat rozumnou
  frekvenci požadavků a identifikovat crawler.

  \item \term{Co je AUC, ROC křivka a confusion matrix?}

  Confusion matrix vznikne porovnáním skutečných tříd s predikovanými třídami při zvoleném prahu.
  ROC křivka zobrazuje TPR proti FPR pro různé prahy. Jeden bod na ROC odpovídá jednomu prahu, ne
  nově trénovanému klasifikátoru. AUC je plocha pod ROC křivkou a vyjadřuje schopnost modelu řadit
  pozitivní případy před negativní.

  \item \term{Co je TF-IDF?}

  TF-IDF je reprezentace textu založená na četnosti termu v dokumentu a jeho vzácnosti v kolekci.
  TF říká, jak často se slovo v dokumentu vyskytuje, IDF snižuje váhu slov běžných napříč dokumenty.
  Výsledkem je řídký vektor dokumentu vhodný pro klasifikaci, vyhledávání nebo podobnost dokumentů.
\end{enumerate}

\section{Zamazal, Tomanová, Vencovský}
Uvedená komise: doc. Ing. Ondřej Zamazal, Ph.D.; Ing. Petra Tomanová, Ph.D., MSc; Ing. Filip Vencovský, Ph.D.

\begin{enumerate}
  \item \term{Jaké architektury neuronových sítí znáte?}

  Základní architektury jsou dopředné plně propojené sítě, konvoluční sítě pro obraz a lokální vzory,
  rekurentní sítě pro sekvence, LSTM/GRU jako stabilnější rekurentní varianty, autoenkodéry, grafové
  neuronové sítě a transformery. Volba závisí na typu dat: tabulky, obraz, text, časové řady, grafy
  nebo sekvence.

  \item \term{Jak fungují rekurentní sítě?}

  Rekurentní síť zpracovává sekvenci postupně a udržuje skrytý stav, který nese informaci z předchozích
  kroků. Hodí se pro text, časové řady a jiné sekvence. Základní RNN trpí problémem mizejících nebo
  explodujících gradientů, proto se používají LSTM a GRU s bránami, které lépe řídí zapamatování a
  zapomínání informace.

  \item \term{Jaké jsou obecné zásady vizualizace dat a prezentace informací?}

  Graf má odpovídat otázce, publiku a typu dat. Je potřeba volit vhodný graf, jasně popsat osy,
  jednotky a zdroj, nepřetěžovat barvami, zvýraznit hlavní sdělení a nezakrývat nejistotu. U časových
  řad je důležité ukázat trend, sezónnost a změny měřítka; u srovnání kategorií jsou často lepší
  sloupcové nebo bodové grafy než koláče.

  \item \term{Jaká rizika mají výstupy LLM při dalším zpracování?}

  LLM může halucinovat, měnit formát odpovědi, generovat nepřesné citace, převzít bias z dat, být
  citlivý na prompt a produkovat věrohodně znějící nepravdy. Při automatickém dalším zpracování je
  nutné validovat strukturu, kontrolovat fakta, používat schémata, testy, lidskou kontrolu u rizikových
  výstupů a uchovávat auditní stopu.

  \item \term{Jak modelovat binární proměnnou TRUE/FALSE regresními metodami?}

  Přirozenou volbou je logistická regrese, která modeluje pravděpodobnost jevu v intervalu 0 až 1.
  Alternativou je probit model nebo lineární pravděpodobnostní model. Lineární model je jednoduchý,
  ale může predikovat hodnoty mimo interval a mívá heteroskedasticitu. Logit a probit jsou vhodnější
  pro interpretaci pravděpodobností a klasifikaci.

  \item \term{Jak obecně probíhá trénování neuronových sítí?}

  Síť provede dopředný průchod, spočítá predikci a ztrátovou funkci. Poté se pomocí backpropagation
  vypočítají gradienty ztráty podle vah a optimalizační algoritmus, například SGD nebo Adam, váhy
  upraví. Proces se opakuje po dávkách dat a epochách. Sleduje se trénovací a validační chyba, aby
  se odhalilo přeučení.

  \item \term{Je rozumné modelovat ceny akcií lineární regresí?}

  Pro krátkodobou predikci cen akcií je jednoduchá lineární regrese často problematická. Ceny bývají
  nestacionární, mají volatilitu měnící se v čase, šum a slabý predikční signál. Smysluplnější je
  pracovat s výnosy než s cenami, testovat stacionaritu, používat validační období v čase a srovnat
  model s jednoduchým benchmarkem.

  \item \term{Co je stacionarita?}

  Stacionární časová řada má stabilní statistické vlastnosti v čase, zejména střední hodnotu, rozptyl
  a autokovariance. Slabá stacionarita znamená konstantní střední hodnotu a rozptyl a kovarianci
  závislou jen na zpoždění. Nestacionaritu lze řešit diferencováním, odstraněním trendu, transformací
  nebo modely, které ji výslovně zohledňují.

  \item \term{Vysvětlete Transformer a self-attention.}

  Transformer zpracovává sekvenci pomocí attention mechanismu místo rekurence. Self-attention počítá,
  jak silně se jednotlivé tokeny mají navzájem vnímat. Tokeny se převádějí na query, key a value
  vektory; podobnost query a key určí váhy, kterými se zprůměrují value vektory. Díky tomu model
  zachycuje dlouhé závislosti a dobře paralelizuje trénování.

  \item \term{Jak zapojit analýzu sentimentu do modelu?}

  Texty lze označit sentimentem ručně, lexikonem nebo klasifikačním modelem. Výstupem může být třída
  pozitivní/neutrální/negativní nebo plynulé skóre. Toto skóre lze použít jako vysvětlující proměnnou,
  agregovat v čase, napojit na události a validovat proti ručně anotovanému vzorku. Je nutné hlídat
  jazyk, doménu, sarkasmus, negaci a posun významu.
\end{enumerate}

\section{Zouhar, Vencovský, Zamazal}
Uvedená komise: doc. Ing. Jan Zouhar, Ph.D.; Ing. Filip Vencovský, Ph.D.; doc. Ing. Ondřej Zamazal, Ph.D.

\begin{enumerate}
  \item \term{Používá se lineární regrese jen pro lineární vztahy?}

  Lineární regrese je lineární v parametrech, ne nutně v původních proměnných. Lze zahrnout logaritmy,
  kvadratické členy, polynomy, spliny, interakce nebo transformace proměnných. Model je stále lineární,
  pokud parametry vstupují lineárně. Například \(y=\beta_0+\beta_1x+\beta_2x^2+u\) je lineární
  regresní model v parametrech.

  \item \term{Jaký je rozdíl mezi lineární a logistickou regresí?}

  Lineární regrese modeluje střední hodnotu spojité závislé proměnné. Logistická regrese modeluje
  pravděpodobnost binárního jevu přes logistickou funkci a log-odds. Lineární regrese minimalizuje
  čtverce chyb, logistická se obvykle odhaduje maximální věrohodností. Logistická regrese je vhodná,
  když je výstup 0/1.

  \item \term{Jaký je vztah mezi lineární regresí a neuronovými sítěmi?}

  Lineární regrese je velmi jednoduchý model: lineární kombinace vstupů bez skrytých vrstev. Neuronová
  síť lze chápat jako skládání lineárních transformací a nelineárních aktivačních funkcí. Bez nelineárních
  aktivací by vícevrstvá síť byla opět jen lineární transformace. Nelinearity dávají sítím schopnost
  aproximovat složité vztahy.

  \item \term{Jaká transakční data jsou vhodná pro predikci odchodu zákazníků?}

  Užitečná jsou data o nákupech, frekvenci používání služby, poslední aktivitě, reklamacích, platbách,
  komunikaci se zákaznickou podporou, změnách tarifu, kampaních, návštěvách webu a smluvních událostech.
  Důležitá je historie v čase, protože churn je často proces postupného slábnutí aktivity. Kvalita,
  aktuálnost a správné časové členění dat zásadně ovlivňují predikci.
\end{enumerate}

\section{Zouhar, Kliegr, Kučera}
Uvedená komise: doc. Ing. Jan Zouhar, Ph.D.; doc. Ing. Tomáš Kliegr, Ph.D.; Ing. Jan Kučera, Ph.D.

\begin{enumerate}
  \item \term{Vysvětlete formáty JSON, XML a CSV.}

  CSV je jednoduchý tabulkový textový formát, kde řádky reprezentují záznamy a sloupce jsou oddělené
  oddělovačem. JSON reprezentuje objekty a pole pomocí dvojic klíč-hodnota a je vhodný pro hierarchická
  data a API. XML používá značky, atributy a stromovou strukturu; je formálnější a často spojovaný se
  schématy, ale bývá objemnější.

  \item \term{Jaká je vnitřní struktura JSON?}

  JSON obsahuje objekty ve složených závorkách, pole v hranatých závorkách, řetězce, čísla, logické
  hodnoty a null. Objekt je množina dvojic klíč-hodnota, kde klíč je řetězec.
  Hodnota může být skalární, objekt nebo pole. JSON je častý pro webová API, protože je čitelný a
  dobře podporovaný programovacími jazyky.

  \item \term{Jaké jsou základní datové struktury například v Pythonu?}

  Mezi základní struktury patří list pro uspořádanou měnitelnou sekvenci, tuple pro neměnitelnou
  sekvenci, dict pro mapování klíč-hodnota, set pro množinu unikátních prvků a string pro text.
  V datové analýze se často používají také tabulkové struktury jako DataFrame a numerická pole.

  \item \term{Jakými metodami zkoumat souvislost a kauzalitu mezi dvěma proměnnými?}

  Pro popis souvislosti lze použít korelaci, kontingenční tabulku, scatterplot, regresi nebo testy
  rozdílů mezi skupinami. Kauzalita vyžaduje silnější design: experiment, přirozený experiment,
  instrumentální proměnné, difference-in-differences, regression discontinuity nebo důkladnou kontrolu
  confounderů. Korelace sama o sobě kauzalitu nedokazuje.

  \item \term{Lze zobecnit výsledek regrese na celou populaci?}

  Ano, pokud jsou splněny předpoklady pro inferenci: vhodný výběr nebo věcně obhajitelný datový
  mechanismus, správná specifikace modelu, exogenita, dostatečná kvalita dat a správně počítané
  standardní chyby. Pokud data nejsou reprezentativní, je zobecnění omezené na populaci podobnou
  pozorovanému vzorku.

  \item \term{Co znamená koeficient $R^2$?}

  $R^2$ vyjadřuje podíl variability závislé proměnné vysvětlený modelem v rámci daného vzorku. Hodnota
  blízká 1 znamená dobrou shodu s pozorovanými daty, ale neznamená kauzalitu ani dobrou predikci mimo
  vzorek. Při přidávání proměnných běžné $R^2$ neklesá, proto se sleduje i adjustované $R^2$.

  \item \term{Jaké jsou nástroje statistické indukce?}

  Statistická indukce používá odhady parametrů, intervaly spolehlivosti, testy hypotéz, p-hodnoty,
  standardní chyby, bootstrap a modelové předpoklady. Cílem je z výběrových dat usuzovat na populaci
  a vyjádřit nejistotu závěru.

  \item \term{Uveďte příklad testu statistické hypotézy a vysvětlete standard error.}

  Příkladem je t-test nulové hypotézy, že regresní koeficient je nula. Testová statistika je odhad
  dělený standardní chybou. Standard error je směrodatná odchylka výběrového rozdělení odhadu; měří,
  jak moc by se odhad měnil při opakovaném výběru. Menší standardní chyba znamená přesnější odhad.

  \item \term{Jaká je role EDA v datové vědě?}

  EDA pomáhá pochopit data, ověřit kvalitu, najít odlehlé hodnoty, rozdělení, vztahy, chybějící
  hodnoty a možné hypotézy. Je důležitá před modelováním, protože špatně pochopená data vedou ke
  špatným proměnným, chybné validaci a mylné interpretaci výsledků.

  \item \term{Jak kvantifikovat asociační pravidla a co znamená support?}

  Asociační pravidla se hodnotí pomocí support, confidence, lift a conviction. Support je relativní
  četnost výskytu položek v pravidle v celém datasetu. Nízký support znamená, že pravidlo může být
  zajímavé, ale týká se malé části dat. Confidence vyjadřuje spolehlivost pravidla a lift sílu vztahu
  oproti nezávislosti.

  \item \term{Jak klasifikovat regresní modely a volit je podle typu úlohy?}

  Regresní modely lze členit na lineární, zobecněné lineární, nelineární, penalizované, robustní,
  panelové, časové a strojově učené regresní modely. Pro spojitou proměnnou často stačí lineární
  regrese, pro binární logit/probit, pro počty Poissonova nebo negativně binomická regrese, pro časová
  data dynamické modely a pro predikci složitých nelinearit ensemble nebo neuronové sítě.

  \item \term{Jak transformovat závislé a nezávislé proměnné v lineární regresi?}

  Často se používají logaritmy, kvadratické členy, interakce, standardizace, dummy proměnné, spliny
  nebo agregace. Logaritmus pomáhá interpretovat procentní změny a snižuje šikmost. Kvadratický člen
  zachycuje zakřivení. Interakce říká, že efekt jedné proměnné závisí na hodnotě jiné proměnné.

  \item \term{Jak vypadá SARIMA model?}

  SARIMA se značí \((p,d,q)(P,D,Q)_s\). První trojice popisuje nesezónní autoregresní řád, diferencování
  a klouzavý průměr. Druhá trojice popisuje sezónní AR, sezónní diferenci a sezónní MA složku.
  Parametr \(s\) je délka sezóny, například 4 pro čtvrtletní data.

  \item \term{Jak fungují neuronové sítě, backpropagation a aktivační funkce?}

  Neuronová síť skládá lineární transformace a nelineární aktivační funkce. Dopředný průchod spočítá
  predikci, ztrátová funkce změří chybu a backpropagation spočítá gradienty. Optimalizátor, například
  SGD nebo Adam, upraví váhy. Aktivační funkce jako ReLU, sigmoid nebo tanh dodávají modelu nelinearitu.
\end{enumerate}

\section{Zouhar, Chudán, Vencovský}
Uvedená komise: doc. Ing. Jan Zouhar, Ph.D.; Ing. David Chudán, Ph.D.; Ing. Filip Vencovský, Ph.D.

\begin{enumerate}
  \item \term{Ovlivní logaritmická transformace následné kódování podle kvantilů?}

  Monotónní transformace, jako je logaritmus pro kladné hodnoty, nemění pořadí pozorování. Pokud se
  tedy kategorie tvoří čistě podle kvantilů stejné proměnné, hranice v pořadí zůstanou stejné. Změní
  se ale vzdálenosti mezi hodnotami a interpretace rozdílů. Pro modelování může být logaritmus užitečný,
  i když samotné kvantilové zařazení nezmění.

  \item \term{Jaké způsoby tokenizace a reprezentace textu znáte?}

  Tokenizace může být po slovech, znacích, n-gramech nebo subword jednotkách jako BPE či WordPiece.
  Reprezentace textu může být bag-of-words, TF-IDF, topic model, word embeddings, contextual embeddings
  z modelů typu BERT nebo vektor dokumentu získaný agregací či transformerem. Volba závisí na velikosti
  dat, jazyku, výpočetních zdrojích a úloze.

  \item \term{Jak funguje TF-IDF?}

  TF-IDF váží term podle četnosti v dokumentu a vzácnosti v celém korpusu. Častá slova v jednom
  dokumentu dostanou vyšší váhu, ale slova častá ve všech dokumentech jsou potlačena. Výsledkem je
  vektor dokumentu, který je vhodný pro klasifikaci, podobnost a vyhledávání, ale nezachycuje slovosled
  ani hlubší význam jako kontextové embeddingy.

  \item \term{Na jakém principu funguje BERT a čím je DistilBERT odlišný?}

  BERT je transformerový jazykový model založený na bidirekcionálním self-attention. Učí se kontextové
  reprezentace tokenů, takže význam slova závisí na okolním textu. DistilBERT je menší a rychlejší
  varianta vzniklá distilací znalostí z BERTu. Má méně parametrů, bývá levnější na inference a často
  si zachová velkou část výkonu původního modelu.

  \item \term{Jaké etické a regulační otázky vznikají při sběru a uchovávání dat?}

  Je potřeba řešit účel zpracování, právní titul, minimalizaci dat, informovaný souhlas nebo jiný
  základ, bezpečnost, omezení doby uchování, práva subjektů údajů, audit, přístupová oprávnění a
  spravedlivé použití modelů. Eticky je důležité, aby data nebyla používána k nečekaným, škodlivým
  nebo diskriminačním účelům.

  \item \term{Jak se liší kvantové ML od klasického deep learningu?}

  Klasický deep learning běží na klasických výpočetních zařízeních a optimalizuje parametrické modely
  pomocí gradientních metod. Kvantové ML využívá kvantové stavy, superpozici a kvantové obvody, často
  v hybridních schématech. Potenciální výhody se týkají některých speciálních výpočtů, ale praktické
  limity jsou šum, malý počet qubitů, obtížná škálovatelnost a nejasná výhoda pro běžné datové úlohy.

  \item \term{Jak se liší ML a DL v možnostech paralelizace?}

  Mnoho klasických ML metod lze paralelizovat přes vzorky, stromy, foldy validace nebo grid search.
  Deep learning je přirozeně paralelizovatelný v maticových operacích na GPU/TPU a přes mini-batche.
  U velkých modelů se používá datová paralelizace, modelová paralelizace a pipeline paralelizace.
  Limitem je komunikace gradientů, paměť a synchronizace.

  \item \term{Jak odhadovat časovou náročnost prolomení algoritmu z časových dat?}

  Záleží na povaze dat. Pokud jde o čas do události, lze použít survival analýzu, hazardní modely nebo
  regresi času do prolomení. Pokud jde o opakované pokusy v čase, lze použít časové řady nebo model
  počtů. Důležité je zachytit cenzorování, velikost instance, parametry algoritmu, hardware a nejistotu
  odhadu.

  \item \term{Podle čeho volit lineární nebo zobecněný lineární model?}

  Rozhoduje typ závislé proměnné a rozdělení chyby. Pro spojitou přibližně normální proměnnou je
  vhodná lineární regrese. Pro binární proměnnou logit/probit, pro počty Poissonova nebo negativně
  binomická regrese, pro podíly a pravděpodobnosti speciální modely podle kontextu. GLM používá link
  funkci a rozdělení z exponenciální rodiny.

  \item \term{Jaké unsupervised metody lze použít pro obrazová data?}

  Lze použít shlukování embeddingů, PCA nebo autoenkodéry pro redukci dimenze, self-supervised učení,
  segmentaci obrazu, detekci anomálií a generativní modely. Často se nejprve obraz převede na vektor
  příznaků pomocí CNN nebo vision transformeru a teprve nad embeddingy se provádí shlukování či
  vyhledávání podobnosti.

  \item \term{Jak funguje AR(1) proces?}

  AR(1) model má tvar
  \[
    y_t = c + \varphi y_{t-1} + \varepsilon_t.
  \]
  Aktuální hodnota závisí na předchozí hodnotě a náhodném šoku. Pokud \(|\varphi|<1\), proces je
  stacionární a šoky postupně odeznívají. Čím je \(\varphi\) blíže 1, tím je řada perzistentnější.

  \item \term{Jaký typ regrese použít pro predikci interakcí na sociálních sítích?}

  Záleží na cílové proměnné. Pro počty reakcí je vhodná Poissonova nebo negativně binomická regrese,
  protože jde o nezáporné počty a data bývají přerozptýlená. Pro binární vysoký/nízký engagement
  logistická regrese. Pro silně nelineární predikci lze použít gradient boosting nebo neuronové sítě.

  \item \term{Jak fungují enkodéry a dekodéry?}

  Enkodér převádí vstup na vnitřní reprezentaci, která má zachytit podstatné informace. Dekodér z této
  reprezentace generuje výstup, například rekonstruovaný vstup, překlad nebo text. V autoenkodéru se
  učí kompaktní latentní reprezentace, v transformerových seq2seq modelech enkodér čte vstupní sekvenci
  a dekodér generuje cílovou sekvenci.

  \item \term{Jak se používají embeddingy v textových úlohách?}

  Embedding převádí slovo, token, větu nebo dokument na hustý vektor. Podobné významy by měly mít
  podobné vektory. Embeddingy se používají pro klasifikaci textů, vyhledávání podobných dokumentů,
  doporučování, clustering, RAG a vstupy do neuronových modelů. Kontextové embeddingy mění reprezentaci
  podle okolního textu.

  \item \term{Proč se zabývat redukcí rozměru dat?}

  Redukce rozměru snižuje počet proměnných při zachování důležité informace. Pomáhá vizualizaci,
  odstranění redundance, snížení šumu, rychlejšímu výpočtu a někdy lepší generalizaci. Typické metody
  jsou PCA, faktorová analýza, autoenkodéry, t-SNE a UMAP. Je ale nutné hlídat ztrátu interpretace.
\end{enumerate}

\section{Zouhar, Chudán, Vencovský: databáze a evaluace}
Tato skupina pochází ze stejné komise jako předchozí část, ale otázky se v protokolech opakovaly
u databází a evaluace modelů, proto jsou pro přehlednost oddělené.

\begin{enumerate}
  \item \term{Jaké typy databází znáte?}

  Základní jsou relační databáze, dokumentové databáze, key-value úložiště, sloupcové databáze,
  grafové databáze, časověřadé databáze, fulltextové indexy a analytické datové sklady. Volba závisí
  na struktuře dat, transakčních požadavcích, dotazech, objemu, latenci a potřebě konzistence.

  \item \term{Popište relační databáze.}

  Relační databáze ukládá data do tabulek s řádky a sloupci. Vztahy se vyjadřují primárními a cizími
  klíči. Důležité jsou integritní omezení, normalizace, SQL, transakce a ACID vlastnosti. Výhodou je
  pevné schéma, konzistence a silná podpora dotazů; nevýhodou může být horší flexibilita pro velmi
  proměnlivá nebo hierarchická data.

  \item \term{Vysvětlete AUC a ROC křivku. Co představuje bod na ROC křivce?}

  ROC křivka vzniká změnou klasifikačního prahu a zobrazuje true positive rate proti false positive
  rate. Jeden bod odpovídá jednomu prahu. AUC je plocha pod ROC křivkou; hodnota 0,5 odpovídá náhodnému
  řazení, 1 dokonalému řazení. AUC je nezávislá na konkrétním prahu, ale nemusí odpovídat business
  nákladům chyb.
\end{enumerate}

\section*{Krátký seznam vynechaných typů otázek}
\addcontentsline{toc}{section}{Krátký seznam vynechaných typů otázek}

Vynechané byly zejména:
\begin{itemize}
  \item otázky z obhajoby závěrečné práce;
  \item dotazy typu \uv{jak byste to použil ve své práci}, pokud bez konkrétní práce neměly obecnou
  odpověď;
  \item jednověté dotazy bez antecedentu, například \uv{uveďte příklad} nebo \uv{co je jeho cílem};
  \item otázky k individuálním výsledkům studenta, konkrétním datasetům, použitým šumům, scénářům
  nebo závěrům práce.
\end{itemize}

\end{document}
